\section{Related work}
\label{sec:related}

\paragraph{Online Bayesian and entropic learning.} The update we use is the
optimised online algorithm for the student--teacher problem, given a Bayesian
reading by \citet{opper1996online,opper1998bayesian} and \citet{solla1999optimal}
and an entropic-dynamics formulation by \citet{caticha2020entropic}; the
identification of the inferred channel-noise level with distrust, and its
promotion to a dynamical variable, is due to \citet{alves2016distrust}. We
inherit the rule and claim nothing about it. What is new here is what it does to
a population when one class-correlated coordinate is added to the
representation, and the analysis of that is ours.

\paragraph{Neural-agent opinion and trust dynamics.} Representing agents as
classifiers over an issue space, rather than as scalars or spins, makes partial
agreement possible \citep{caticha2011agent,caticha2015whom}, which is what lets
one ask whether distrust follows disagreement or the reverse. Earlier work in
this family studies societies without any group label
\citep{simoes2018mean,caticha2019trust}. The label, the representation error that
carries it, and the collective consequences are the subject here.

\paragraph{Opinion dynamics and structural balance.} Models of consensus and
polarisation run from bounded-confidence rules
\citep{deffuant2000mixing,hegselmann2002opinion} through cultural dissemination
\citep{axelrod1997dissemination} to reinforcement-driven echo chambers
\citep{baumann2020modeling}; see \citet{castellano2009statistical} for a survey.
Our triple-level diagnostics are those of structural balance
\citep{heider1946attitudes,cartwright1956structural,antal2005dynamics,marvel2011continuous},
but with one difference worth stating: there the sign structure is the primitive
object, whereas here both sign structures are \emph{generated} by inference. We
also study a transient at a fixed horizon rather than an equilibrium, because the
dynamics anneals and has no stationary state.

\paragraph{Fairness, sensitive attributes and spurious features.} The dominant
framing locates group bias in data and in single-model decisions, as
representational harm \citep{blodgett2020language,bender2021dangers} or as
allocative unfairness with formal criteria attached
\citep{dwork2012fairness,hardt2016equality,barocas2023fairness}. Our mechanism is
complementary: the outcome is produced by an irrelevant feature interacting with
the \emph{learning rule} across a population, and it is invisible to an audit
that inspects agents one at a time. That irrelevant features degrade a classifier
is old; that they can restructure a society of classifiers is the point here.

\paragraph{Collective effects of local rules.} The general phenomenon --- locally
sensible optimisation producing undesirable system-level behaviour --- is
recognised across multi-agent learning, but we are not aware of a treatment in
which the coupling runs through an \emph{inferred} trust between learners. That
is the gap this paper occupies.
